% =====================================================================
%  c-refs.tex —— 参考文献（底本第 7 章，英文原貌；由 tools/gen_refs.py 生成）
%  说明：按学术惯例不翻译；仅做 LaTeX 转义与重排，未逐条校对（见 README 已知限制）。
%  共 32 条。
% =====================================================================
\section*{参考文献}
\addcontentsline{toc}{section}{参考文献}

\begingroup
\small
\sloppy

\refitem{Abo11}{Abouzaid, M. (2011). A cotangent fibre generates the Fukaya category. Advances in Mathematics, 228(2), 894--939. doi:10.1016/j.aim.2011.06.007}
\refitem{Abo12}{Abouzaid, M. (2012). On the wrapped Fukaya category and based loops. The Journal of Symplectic Geometry, 10(1), 27--79. doi:10.4310/jsg.2012.v10.n1.a3}
\refitem{Abo14}{Abouzaid, M. (2014). Family Floer cohomology and mirror symmetry. Proceedings of the International Congress of Mathematicians---Seoul 2014. Vol. II, 813--836P. Retrieved from http://www.ams.org/mathscinet-getitem?mr=3728639}
\refitem{Abo21}{Abouzaid, M. (2021). Homological mirror symmetry without correction. Journal of the American Mathematical Society, 34(4), 1059--1173. doi:10.1090/jams/973}
\refitem{Ati82}{Atiyah, M. F. (1982). Convexity and commuting Hamiltonians. The Bulletin of the London Mathematical Society, 14(1), 1--15. doi:10.1112/blms/14.1.1}
\refitem{AEK16}{Abouzaid, M., Auroux, D., \& Katzarkov, L. (2016). Lagrangian fibrations on blowups of toric varieties and mirror symmetry for hypersurfaces. Publications Mathématiques. Institut de Hautes Études Scientifiques, 123, 199--282. doi:10.1007/s10240-016-0081-9}
\refitem{AS21}{Abouzaid, M., \& Sylvan, Z. (2021). Homological mirror symmetry for local SYZ singularities. arXiv. Retrieved from http://arxiv.org/abs/2107.05068}
\refitem{Aur07}{Auroux, D. (2007). Mirror symmetry and T-duality in the complement of an anticanonical divisor. \emph{Journal of Gökova Geometry Topology}, \emph{1}, 51--91. Retrieved from http://www.ams.org/mathscinet-getitem?mr=2386535}
\refitem{Aur25}{Auroux, D. (2025). Holomorphic discs of negative Maslov index and extended deformations in mirror symmetry. arXiv. Retrieved from http://arxiv.org/abs/2309.13010}
\refitem{BO01}{Bondal, A., \& Orlov, D. (2001). Reconstruction of a variety from the derived category and groups of autoequivalences. \emph{Compositio Mathematica}, \emph{125}(3), 327--344. doi:10.1023/A:1002470302976}
\refitem{CO06}{Cho, C.-H., \& Oh, Y.-G. (2006). Floer cohomology and disc instantons of Lagrangian torus fibers in Fano toric manifolds. \emph{Asian Journal of Mathematics}, \emph{10}(4), 773--814. doi:10.4310/AJM.2006.v10.n4.a10}
\refitem{Eva23}{Evans, J. (2023). \emph{Lectures on Lagrangian torus fibrations} (p. xiii+225). doi:10.1017/9781009372671}
\refitem{Flo89}{Floer, A. (1989). Witten's complex and infinite-dimensional Morse theory. \emph{Journal of Differential Geometry}, \emph{30}(1), 207--221. Retrieved from http://projecteuclid.org/euclid.jdg/1214443291}
\refitem{FLTZ11}{Fang, B., Liu, C.-C. M., Treumann, D., \& Zaslow, E. (2011). A categorification of Morelli's theorem. \emph{Inventiones Mathematicae}, \emph{186}(1), 79--114. doi:10.1007/s00222-011-0315-x}
\refitem{FOOO09}{Fukaya, K., Oh, Y.-G., Ohta, H., \& Ono, K. (2009). \emph{Lagrangian intersection Floer theory: anomaly and obstruction. Part I} (p. xii+396). doi:10.1090/amsip/046.1}
\refitem{Fuk09}{Fukaya, K. (2009). Lagrangian surgery and rigid analytic family of Floer homologies. Retrieved from https://www.math.kyoto-u.ac.jp/\textasciitilde{}fukaya/Berkeley.pdf}
\refitem{GPS24}{Ganatra, S., Pardon, J., \& Shende, V. (2024). Sectorial descent for wrapped Fukaya categories. Journal of the American Mathematical Society, 37(2), 499--635. doi:10.1090/jams/1035}
\refitem{GS82}{Guillemin, V., \& Sternberg, S. (1982). Convexity properties of the moment mapping. Inventiones Mathematicae, 67(3), 491--513. doi:10.1007/BF01398933}
\refitem{Jin25}{Jing, H. (2025). On SYZ mirrors of Hirzebruch surfaces. arXiv. Retrieved from http://arxiv.org/abs/2504.01889}
\refitem{Kon94}{Kontsevich, M. (1994). Homological algebra of mirror symmetry. arXiv. Retrieved from https://arxiv.org/abs/alg-geom/9411018v1}
\refitem{KS06}{Kontsevich, M., \& Soibelman, Y. (2006). Affine structures and non-Archimedean analytic spaces. In \emph{The unity of mathematics} (pp. 321--385). doi:10.1007/0-8176-4467-9\_9}
\refitem{Kuw20}{Kuwagaki, T. (2020). The nonequivariant coherent-constructible correspondence for toric stacks. \emph{Duke Mathematical Journal}, \emph{169}(11), 2125--2197. doi:10.1215/00127094-2020-0011}
\refitem{LS03}{Leung, N. C., \& Symington, M. (2010). Almost toric symplectic four-manifolds. \emph{The Journal of Symplectic Geometry}, \emph{8}(2), 143--187. doi:10.4310/jsg.2010.v8.n2.a2}
\refitem{LZ20}{Lau, S.-C., \& Zheng, X. (2020). SYZ mirror symmetry for hypertoric varieties. \emph{Communications in Mathematical Physics}, \emph{373}(3), 1133--1166. doi:10.1007/s00220-019-03535-z}
\refitem{MS04}{McDuff, D., \& Salamon, D. (2004). J-holomorphic curves and symplectic topology (p. xii+669). doi:10.1090/coll/052}
\refitem{ND20}{nowhere dense. (2020). The Zariski topology of a variety is determined by its principal Cartier divisors. MathOverflow. Retrieved from https://mathoverflow.net/q/355978}
\refitem{Ngoc03}{Vũ Ngọc, S. (2003). On semi-global invariants for focus-focus singularities. Topology. An International Journal of Mathematics, 42(2), 365--380. doi:10.1016/S0040-9383(01)00026-X}
\refitem{Oh96}{Oh, Y.-G. (1996). Floer cohomology, spectral sequences, and the Maslov class of Lagrangian embeddings. \emph{International Mathematics Research Notices}, \emph{1996}(7), 305--346. doi:10.1155/S1073792896000219}
\refitem{PZ98}{Polishchuk, A., \& Zaslow, E. (1998). Categorical mirror symmetry: The elliptic curve. \emph{Advances in Theoretical and Mathematical Physics}, \emph{2}(2), 443--470. doi:10.4310/ATMP.1998.v2.n2.a9}
\refitem{SYZ96}{Strominger, A., Yau, S. T., \& Zaslow, E. (1996). Mirror symmetry is T-duality. \emph{Nuclear Physics B}, \emph{479}(1-2), 243--259.}
\refitem{Sym03}{Symington, M. (2003). Four dimensions from two in symplectic topology. In \emph{Topology and geometry of manifolds (Athens, GA, 2001)} (pp. 153--208). doi:10.1090/pspum/071/2024634}
\refitem{TY02}{Thomas, R. P., \& Yau, S.-T. (2002). Special Lagrangians, stable bundles and mean curvature flow. \emph{Communications in Analysis and Geometry}, \emph{10}(5), 1075--1113. doi:10.4310/CAG.2002.v10.n5.a8}

\endgroup
